Afirma-se com frequência que pipelines de dados escritos como composição de funções puras são seguros para reexecução, mas essa afirmação mistura práticas de natureza distinta. Este artigo investiga quais práticas associadas à programação funcional determinam a idempotência de um pipeline em lote e quais decorrem de decisões de engenharia independentes de paradigma. Um pipeline de eventos de pedidos em Python foi parametrizado por quatro fatores binários (tempo injetado, transformações puras, escrita por sobrescrita de partição e deduplicação por chave) e as 16 combinações foram submetidas a seis cenários de execução, com 30 conjuntos de dados sintéticos cada e um gabarito calculado de forma independente. As hipóteses foram registradas antes das medições. Nas 2.880 execuções, os resultados coincidiram com a previsão em todas as 96 células. Nenhum dos dois grupos de práticas basta sozinho: as práticas funcionais foram necessárias em três dos cinco cenários de reexecução, as de engenharia em quatro, e ambas em dois. A necessidade das transformações puras depende de o estado em memória sobreviver entre execuções. A avaliação de testabilidade está em andamento.
